\documentclass[12pt,a4paper,oneside]{report}

% Codificación y idioma
\usepackage[T1]{fontenc}
\usepackage[utf8]{inputenc}
\usepackage[spanish]{babel}

% Márgenes y microtipografía
\usepackage[a4paper,margin=2.5cm]{geometry}
\usepackage{microtype}

% Gráficos, tablas y matemáticas
\usepackage{graphicx}
\usepackage{caption}
\usepackage{subcaption}
\usepackage{booktabs}
\usepackage{tabularx}
\usepackage{amsmath,amssymb}

% Hipervínculos
\usepackage[hidelinks]{hyperref}

% Numeración de secciones sin dependencia de capítulos (1,2,3... en vez de 0.1,0.2)
\usepackage{chngcntr}
\counterwithout{section}{chapter}

\begin{document}


\begin{table}[ht]
    % Portada en la misma primera página (sin página separada)
\begingroup
\begin{center}
	{\LARGE \textbf{Cronograma final del TFG}\par}
	\vspace{1em}
	{\large Raúl Martín-Romo Sánchez\par}
    {\noindent\textbf{Entrega:} última semana de enero}
\end{center}
\bigskip
\endgroup

  \centering
  \small
  \begin{tabularx}{\textwidth}{@{}p{4.5cm} p{3.2cm} X@{}}
    \toprule
    	\textbf{Fase} & \textbf{Fechas} & \textbf{Actividades} \\
    \midrule

    1. Reuniones iniciales y elección del tema
      & 6 - 14 de noviembre
      & \begin{minipage}[t]{\linewidth}
          \begin{itemize}
            \item 6/11: Primera reunión con los guías caninos. Discusión de líneas de trabajo.
            \item 14/11: Elección definitiva del tema del TFG.\end{itemize}
        \end{minipage} \\

    \midrule

    2. Revisión bibliográfica y estado del arte
      & 20 de noviembre - 4 de diciembre
      & \begin{minipage}[t]{\linewidth}
          \begin{itemize}
            \item Revisión de literatura sobre narices electrónicas.
            \item Revisión de trabajos de Jesús Lozano Rogado.
            \item Revisión sobre detección canina.
            \item Redacción del marco teórico.
          \end{itemize}
        \end{minipage} \\

    \midrule

    3. Formación y entrenamiento de la nariz electrónica
      & 5 de diciembre - 10 de enero
      & \begin{minipage}[t]{\linewidth}
          \begin{itemize}
            \item Reunión con Jesús Lozano Rogado.
            \item Formación con Jesús Lozano Rogado.
            \item Desarrollo de la nariz electrónica.
            \item Obtención de patrones de olor.
            \item Entrenamiento del modelo.
          \end{itemize}
        \end{minipage} \\

    \midrule

    4. Diseño del enfoque experimental, comparación y analisis de datos
      & 11 - 18 de enero
      & \begin{minipage}[t]{\linewidth}
          \begin{itemize}
            \item Definición final de la metodología comparativa.
            \item Coordinación con guías caninos.
            \item Pruebas comparativas (nariz electrónica vs perros) y recogida de datos.
            \item Limpieza y análisis de datos.
            \item Comparación entre ambos sistemas.
            \item Elaboración de gráficas y tablas.
          \end{itemize}
        \end{minipage} \\

    \midrule

    5. Redacción final del TFG
      & 19 - 31 de enero
      & \begin{minipage}[t]{\linewidth}
          \begin{itemize}
            \item Redacción final.
            \item Integración completa del documento.
            \item Correcciones finales.
          \end{itemize}
        \end{minipage} \\

    \midrule

    6. Preparación de la presentación
      & Hasta el dia de la presentación
      & \begin{minipage}[t]{\linewidth}
          \begin{itemize}
            \item Diseño de las diapositivas para la presentación.
            \item Preparación de la exposición oral (15--20 minutos).
          \end{itemize}
        \end{minipage} \\

    \bottomrule
    \end{tabularx}
    \end{table}

\end{document}